% Nouveau EGA V (cote 25): archive cote 25, site page 152 (PDF page 153).
% typescript. Draft compared: none.
% AI-transcribed by Claude Opus 5.5 (claude-opus-5-5, adaptive thinking, effort high), Message Batches API, batch ega25b (msgbatch_01A7n5YuXubUPNK1zsepi8RG), 2026-11-13; not yet checked by eye.
\scanpage{25/152}
\mspage{13}

% NOTE: a long handwritten diagonal stroke runs from the top edge (above "b) l'image") down to the left margin just above "Notes que", across the whole passage from "soit" to "non dégénérée."; it is transcribed as a deletion of that passage
\deleted{\inserted[typed]{= $A_{m-1}$}}
\deleted{soit \xout{\unsure{singulier}} régulier de dimension $n-(m-1)$, et b) l'image de $u(\phi_m)$}
\deleted{dans $A_{m-1}$ admette le point fermé de Spec($A_{m-1}$) comme \uline{zéro singulier}}
\deleted{\uline{ordinaire} (cf définition 1).}

\deleted{(ii bis) \xout{Pour toute} Le rang de Im($E\otimes k \to V$) est m-1, et pour}
\deleted{toute base $(\phi_i)_{1\leq i\leq m}$ de E telle que les images des $\phi_i$ $(1\leq i\leq m-1)$ \inserted[typed]{dans V}}
\deleted{sont linéairement indépendantes,\xout{\unsure{e}} les conditions a) et b) de (ii) sont}
\deleted{satisfaites.}

\deleted{(iii) Le noyau \xout{de} K de $E\otimes k \to V$ est de dimension 1, et si $\phi\in K$\hand{$-\{0\}$,}}
\deleted{\xout{la forme} l'élément de Symm$^2$(W) défini par $\phi$ est une forme quadratique}
\deleted{non dégénérée.}

% NOTE: a handwritten double vertical line in the left margin runs from "Notes que" down to "= dim O_{X,x}", with a hook-shaped mark at its top
Notes que l'analogue des conditions (ii) et (ii bis) aurait déjà
pu remonter dans la proposition 9. Il y a lieu d'introduire la terminolo-%
gie suivante: si X est un préschéma régulier, $\underline{E}$ un Module localement libre
\xout{de rang m} sur X, u: $\underline{E}\to\underline{O}_X$ un homomorphisme i.e. u une section de
$\underline{E}$\hand{$^\vee$} , on dira que l'élément x de $V(u)=V(u(\underline{E})\underline{O}_X)$ est un \xout{zéro non sing}
``\uline{zéro non singulier}'' pour u\hand{,} si u est $\underline{O}_X$-régulier i.e. dim $\underline{O}_{V(u),x}$
= dim $\underline{O}_{X,x}$ - rang$_x$ $\underline{E}$ , et si de plus x est un point régulier de V(u),
ce qui en termes d'une base $\phi_1,\ldots,\phi_m$ de $\underline{E}$ au voisinage de x s'exprime
aussi en disant que les $u(\phi_i)$ forment une partie d'un systême régulier
de paramètres de $\underline{O}_{X,x}$ ; dans le cas contraire, x est dit un \uline{zéro}
\uline{singulier}. On dit que x est un \uline{zéro singulier ordinaire} s'il satisfait
aux conditions équivalentes du corollaire 10 (ce qui implique que c'est
bien un zéro singulier), où évidemment $A=\underline{O}_{X,x}$ etc. On dit encore que
x est un \uline{zéro supersingulier} s'il est singulier, sans \hand{\^{e}}tre singulier
ordinaire. Donc non supersingulier = non singulier ou singulier ordinaire
Enfin, lorsque X est lisse sur un corps k, on introduit les variantes
géométriques de ces notions de façon évidente. On notera que dans ces

\mspage{152}
